\documentclass[../main.tex]{subfiles}
\begin{document}
%%%%%%%%%%%%%%%%%%%%%%%%%%%%%%%%%%%%%%%%%%%%%%%%%%%%%%%%%%%%%%%%%%%%%%%%%%%%%%%%%%%%%%%%%%%%%%%%%%%
\chapter{Using OmniGraffle in Presentation Mode (\texorpdfstring{\textbf{PRO}}{PRO})}

After creating your design---whether it's a prototype for an app or website, or a series of
posters to promote an upcoming event---chances are you'll want to present that design to
management or to a client. Fortunately for you, we have built in \textbf{Presentation Mode} to
OmniGraffle.

This chapter explains how to get the most from Presentation Mode, and provides you with tips and
tricks for presenting from within OmniGraffle.

\section{Presentation Basics}

When you enter Presentation Mode, either by choosing \textbf{View} $\blacktriangleright$
\textbf{Start Presentation}, or by pressing \textbf{Option-Command-P}, OmniGraffle takes over
your entire screen. While in Presentation Mode:

\begin{itemize}
  \item The \textbf{Action Browse tool} \ogiconw{19.4pt}{res/og7_tools_actionbrowse.png} is
    sticky; you cannot choose another tool in the Toolbar.
  \item The Toolbar and Inspectors bar are inaccessible.
  \item The Sidebar and menu bar are hidden from sight but remain available if you move your
    mouse to the left or top edge of the screen, respectively.
\end{itemize}

Each canvas in your OmniGraffle document is a \textit{slide} in your presentation. To navigate
through your presentation, you'll use the arrow keys on your keyboard as follows:

\begin{itemize}
  \item Use the Up ($\uparrow$) and Down ($\downarrow$) Arrow keys on your keyboard to go back
    or forward through the canvases in your document, respectively.
  \item Use the Left ($\leftarrow$) and Right ($\rightarrow$) Arrow keys to select the previous
    or next artboard on the current canvas, respectively.

    \begin{itemize}
      \item If you press the Left or Right Arrow keys while on a canvas that does not have an
        artboard, the presentation goes backward or forward through the canvases, respectively, to
        find and select an artboard.
      \item If your document does not contain artboards, nothing happens when you press the Left
        or Right Arrow keys.
    \end{itemize}

    When you select an artboard, the objects on that artboard zoom in to fill the entire screen.
    To zoom back out to see everything on the canvas, press the Up ($\uparrow$) Arrow key, or
    press the Down ($\downarrow$) key to proceed to the next canvas/slide.
\end{itemize}

Move your mouse to the left side of the screen to bring in the Sidebar. The Sidebar only reveals
canvases and artboards while presenting; you cannot access individual layers or objects in the
Sidebar.

Not all menu options are available while in Presentation Mode. For example, you cannot view the
rulers or set alignment guides while presenting. The menu bar goes back into hiding when you
move the mouse away.

To exit Presentation Mode, you can either press \textbf{Option-Command-P} or \textbf{Escape}.

\subsection{Using the Zoom Commands in Presentation Mode}

While presenting, you can use the following keyboard shortcuts for adjusting the magnification
level of objects onscreen:

\begin{itemize}
  \item Zoom to Actual Size (\textbf{Option-Command-0})
  \item Zoom In (\textbf{Shift-Command->})
  \item Zoom Out (\textbf{Shift-Command-<})
\end{itemize}

\section{Configuring OmniGraffle's Presentation Preferences}

To customize the settings for Presentation Mode, you'll use the Presentation preferences,
accessed by choosing \textbf{OmniGraffle} $\blacktriangleright$ \textbf{Preferences}, and then
choose the \textbf{Presentation} tab.

You can choose one of the following options to specify whether objects on the canvas receive a
highlight (border) around the object:

\begin{itemize}
  \item \textbf{Highlight objects by pointing} --- Places a colored highlight around objects
    that you mouse over.
  \item \textbf{Highlight objects by clicking} --- Places a colored highlight around objects
    that you click.
  \item \textbf{Don't highlight objects} --- Objects on the canvas do not receive a highlight.
\end{itemize}

You can adjust the width of the highlight border using the \textbf{Width} slider.

\ogfigure{0.800}{res/og7_present01.png}

Beneath the \textbf{Highlight Behavior} options, the \textbf{Show link badge when highlighted}
preference has been turned on by default. This option applies to objects that are assigned an
Action using the Action inspector. If an object has an action assigned, its Action badge is
displayed in Presentation Mode.

\ogfigure{0.800}{res/og7_present02.png}

You can change the default highlight color by clicking
\ogicon{12.5pt}{res/og7_prefs_highlightwidget.png} to open the Colors window. You can change
the color to anything you'd like, or you can set it so that the highlight surrounding an
object always matches macOS's highlight color (set in \textbf{System Preferences}
$\blacktriangleright$ \textbf{General} $\blacktriangleright$ \textbf{Highlight color}); for
example:

\begin{enumerate}
  \item Under \textbf{Highlight Appearance}, click
    \ogicon{12.5pt}{res/og7_prefs_highlightwidget.png} to open the Colors window.
  \item At the top of the Colors window, click the third tab, \textbf{Color Palettes}.
  \item In the pop-up menu, choose \textbf{Developer}.
  \item Scroll down in the list of color options and choose \textbf{selectedControlColor}.
\end{enumerate}

\ogfigure{0.673}{res/og7_prefs_presentation_highlight.png}

Now whenever you give a presentation from within OmniGraffle---such as when you're demoing an
app prototype---the objects you hover over are highlighted with your system's default
selection color.
%%%%%%%%%%%%%%%%%%%%%%%%%%%%%%%%%%%%%%%%%%%%%%%%%%%%%%%%%%%%%%%%%%%%%%%%%%%%%%%%%%%%%%%%%%%%%%%%%%%
\end{document}
